\documentclass[12pt,a4paper]{article}
\usepackage[UTF8]{ctex}
\usepackage{amsmath,amssymb}
\usepackage{booktabs}
\usepackage{float}
\usepackage{fancyhdr}   % 页眉页脚控制
\usepackage{geometry}
\usepackage{subfigure}  % 子图排版
\usepackage{subcaption}
\usepackage{graphicx}
\usepackage{listings} % 代码排版宏包
\usepackage{xcolor} % 代码高亮颜色
\usepackage{enumitem} % 列表格式控制宏包
\geometry{left=2.5cm,right=2.5cm,top=2.5cm,bottom=2.5cm}

\pagestyle{fancy}
\fancyhf{} % 清空默认页眉页脚
\fancyhead[L]{\small 示波器的使用}       % 左上角页眉
\fancyhead[R]{\small 2513318}          % 右上角章节名
\fancyfoot[C]{\thepage}                   % 居中页码（页脚）
\renewcommand{\headrulewidth}{0.4pt}      % 页眉线粗细
\renewcommand{\footrulewidth}{0.4pt}      % 页脚线粗细
%报告开始
\begin{document}
	%设置课程标题
	\begin{center}
		\heiti\LARGE{《大学基础物理实验》课程实验报告}
	\end{center}
	%设置实验人信息以及实验时间表格
	\begin{center}
		\begin{tabular}{lcr}
			{\songti 姓名及学号：陈秋实2513318}  \quad 专业：密码科学与技术 \quad 组别：L组1号 \\
			{\songti  学院：密码与网络空间安全学院 \quad 实验时间：2026年4月10日~星期五~下午}\\
		\end{tabular}
	\end{center}
	\vspace{-0.2cm}
	{\noindent}	 \rule[-10pt]{16cm}{0.05em}\\
	\vspace{-0.4cm}
	
	%实验题目
	\begin{center}
		\LARGE\textbf{实验题目：示波器的使用}
	\end{center}
	
	%实验原理
	\section{仪器与用具}
	\par 1.1仪器品牌与型号：
	\par 示波器：普源DS1102E，信号发生器：F05函数发生器
	\par 1.2电阻阻值：1k$\Omega$，电容值：0.1$\mu F$
	
	\section{基本使用}
	\par 将信号源($1KHz,3Vp-p$)和变压器电压同时输出到示波器，分别稳定并显示适当的波形。重点熟悉触发对波形的作用。
	
	\section{实验数据}
	\subsubsection*{信号源和变压器测量}
	\par 实验图片如下:
	\begin{figure}[htbp]
		\centering
		\includegraphics[width=0.65\textwidth]{1.jpg}
	\end{figure}
	
	\subsubsection*{李萨如图形测量法}
	\begin{table}[!htbp]
		\centering 
		\begin{tabular}{|c|c|c|c|c|c|}
			\hline
			\rule{0pt}{20pt}$\frac{\mbox{与水平线的交点数}}{\mbox{与竖直线交点数}}=\frac{n_2}{n_1}$    & 1 & 1/2 & 1/3 & 2 & 2/3  \\ \hline
			\rule{0pt}{20pt}函数发生器频率 & 50.01 & 99.94 & 149.91 & 24.98&74.97   \\ \hline
			\rule{0pt}{20pt}算出的市电频率 & 50.01& 49.97 & 49.97 &49.96&49.98 \\ \hline
		\end{tabular}
	\end{table}
	\par 计算得到平均市电频率为49.98Hz
	
	\subsection{测量RC电路的相位差}
	\par 连接电路，将信号发生器频率设定为$f=1.59KHz$
	\begin{figure}[htbp]
		\centering
		\includegraphics[width=0.65\textwidth]{2.jpg}
	\end{figure}
	\subsubsection*{椭圆法}
	$$
	|\theta|=\arcsin\frac{2x_0}{2x_m}=\arcsin\frac{2.88}{3.90}\approx0.738
	$$
	\subsubsection*{位移法}
	\[
	\theta = \frac{l}{l_0}\times 360^o = \frac{224}{312}\approx0.718
	\]
	
	\subsection{思考题}
	$\tilde{u}_R = \frac{R}{R + R_C} \, u = \frac{R}{R + \frac{1}{j\omega C}} \, u,\ \text{其分压值为 } |\tilde{u}_R| = \frac{R}{\sqrt{R^2 + \left( \frac{1}{\omega C} \right)^2}} \, u$
	
	
\end{document}